\section{The remaining distance to quasi-polynomial recognition}
\label{sec:complexity}

\subsection{Three separate obstacles}

A faster exact scalar filter addresses only one part of the recognition
problem. Three distinct missing implications must not be conflated.

First, equality of a selected Jones value to $1$ need not decide the knot.
At $q=5$ this is disproved by Theorem~\ref{thm:blind}; at $q=6$ we assume
no global detection statement. The completeness of the enclosing recognizer
comes from its retained exact fallback, not from these evaluations.

Second, the input diagram does not come with a promise of logarithmic or
polylogarithmic width. Theorems~\ref{thm:transfer} and
\ref{thm:componentwidth} are bounds in a supplied order's measured
parameters. A polynomial heuristic finds useful orders but has no universal
$O((\log n)^2)$ guarantee. Planar separator methods and tensor contraction
are relevant, established tools \cite{maria2021,seymour1994}; using them does
not imply a polylogarithmic bound for arbitrary diagrams.

Third, the full categorical complex contains more information than the
scalar partition function. Boundary size does not alone control object
multiplicities, ranks, differential supports or the size of an explicit
homology basis. Kelomäki and Schütz prove exponential behavior for explicit
scanning or divide-and-conquer constructions even on certain positive and
alternating three-braid families, while giving polynomial algorithms for
compressed three-braid homology and fixed homological windows at fixed strand
count \cite{kelomaki2026}. This is a representation obstruction, not a reason
to abandon compression.

\subsection{A finite exact menu decides Jones-polynomial triviality}

The fixed-color obstruction can be extended to an unconditional decision
about the entire Jones polynomial, although this still does not settle
unknot recognition.

\begin{proposition}[A growing exact specialization menu]\label{prop:menu}
For a knot diagram with $n>0$ crossings, suppose the exact quadratic
evaluations equal $1$ for every integer
$q\in\{5,6,\ldots,2n+5\}$. Then $V_D(t)$ is identically $1$.
Conversely, the constant polynomial $1$ passes every evaluation.
On a supplied order of width $w$, this decision can be carried out using
the one-tensor algorithm in
\[
 n^{\lfloor w/2\rfloor+O(1)}
\]
bit operations. In particular it is quasi-polynomial when
$w=O(\log n)$.
\end{proposition}
\begin{proof}
In a normalized smoothing summand, the $A$ exponents come from a smoothing
weight in $[-n,n]$, the writhe correction in $[-3n,3n]$, and
$\delta^{c-1}$ with $c-1\le n$. Hence all $A$ exponents lie in $[-6n,6n]$.
For a knot they are multiples of four, so all Jones $t$ exponents lie in
the conservative interval $[-2n,2n]$.
Thus $P(t)=t^{2n}(V_D(t)-1)$ is a polynomial of degree at most $4n$.

Equality at a root of $t^2-(q-2)t+1$ makes this irreducible quadratic divide
$P$ over $\Q$. The $2n+1$ quadratics for the specified distinct values of
$q$ are pairwise coprime. If $P\ne0$, their product of degree $4n+2$
cannot divide it. Thus $P=0$.

There are $O(n)$ evaluations with $q=O(n)$. For $f\le w/2\le n$,
the number of equality patterns is at most the Bell number
$B(f)\le f^f\le n^f$, with the empty pattern treated separately.
Introducing two vertices has only $O((f+2)^2)$ canonical branches,
independent of the multiplicities attached to those branches.
The coordinate lengths are $O(n\log n)$ by Theorem~\ref{thm:bits}.
All remaining factors, including exact arithmetic and normalization, are
polynomial in $n$ and $w$.
\end{proof}

A smaller menu follows from the same proof without new topology.
For writhe $r$, the actual conservative exponent interval is
$[\lceil3(r-n)/4\rceil,\lfloor3(r+n)/4\rfloor]$.
It contains zero and has width $D\le\lfloor3n/2\rfloor$.
Thus $\lfloor D/2\rfloor+1$ distinct quadratic evaluations suffice.
The larger $2n+1$ menu keeps the statement independent of this refinement.

This is elementary degree counting, not a new claim about the difficulty of
the Jones polynomial. It shows exactly what a sufficiently large menu can
certify. The menu is a mathematical extension; the default recognizer does
not automatically run $O(n)$ different color counts. The complete
unknot-detection question for $V_D(t)\equiv1$ remains outside its conclusion.
It also shows why a growing number of colors cannot be hidden in a
fixed-$q$ complexity statement.

\subsection{A precise conditional algebraic route}

Reduced Khovanov rank one characterizes the unknot
\cite{kronheimer2011}. The same rank-one decision is sound over $\F_2$:
rank one over $\F_2$ forces rational rank at most one by universal
coefficients, while reduced knot homology is nonzero; the rational
rank-one theorem then applies. Conversely the unknot has reduced rank one
over $\F_2$.
The baseline's local-complex and cancellation approach follows the practical
framework of Bar-Natan~\cite{barnatan2007}; a general small-complex bound
is a separate requirement.

To turn that endpoint into a quasi-polynomial algorithm, a compression
scheme needs a closed set of operations and a size theorem, not merely
small examples or an existence statement.

\begin{theorem}[Conditional summary criterion]\label{thm:conditional}
Let $Q(n)=\exp((\log(n+2))^c)$ for a fixed constant $c$.
Suppose a proposed exact recognition method has, on every $n$-crossing
input, the following proved properties:
\begin{enumerate}
\item it constructs at most $Q(n)$ summaries or intermediate states;
\item each summary has at most $Q(n)$ bits;
\item constructing, composing, restricting, simplifying, comparing, and
checking a summary uses at most $Q(n)^{O(1)}$ bit operations;
\item a final predicate, evaluable within the same bound, is equivalent to
reduced Khovanov rank one or to another complete unknot criterion.
\end{enumerate}
Then the resulting exact recognizer is quasi-polynomial time.
\end{theorem}
\begin{proof}
Multiply the number of operations by their cost. A fixed power of
$\exp((\log(n+2))^c)$ is still quasi-polynomial. Exact invariants and the
terminal equivalence imply correctness.
\end{proof}

The theorem is intentionally a contract rather than a breakthrough: its
hypotheses identify the work still missing. The Potts representations satisfy
analogous storage and operation bounds at bounded width, but their terminal
predicate is incomplete. The residue-polynomial theorem shortens a local
certificate, but does not prove the global summary count or size. The earlier
scalar splitter finds some endomorphisms but does not guarantee a useful
decomposition on every input.

For an explicit ring-valued implementation, the coefficient space alone has
$p^b$ basis monomials. A block with multiplicity $r$ can have $r^2$
entries. Bounding an inverse polynomial's degree by $pr$ leaves both
$p^b$ and the total number of blocks to be controlled.
Using a circuit or shared DAG can reduce their representation, but then
composition, equality and zero detection for that representation also require
proved bounds. Constant-size formulas whose expansion is exponential do not
by themselves satisfy the criterion.

\subsection{The geometric route and the current literature}

Lackenby's 2021 slides announce a deterministic recognition algorithm with
bound $2^{O((\log n)^3)}=n^{O((\log n)^2)}$ \cite{lackenby2021}.
That is quasi-polynomial, but it is not the stronger
$n^{O(\log n)}$ target sometimes used in project notes.
The July 2026 hierarchy preprint supplies a new recognition algorithm and
certification framework; its Section 9 separates iteration counting from
the cost of operations and bounds iterations by $L(G+1)^L$,
where $L$ is maximum hierarchy length and $G$ maximum surface
pattern-complexity \cite{lackenby2026}. We use $G$ here to distinguish it
from the component frontier $g$.

Published compressed-curve algorithms compute intersection numbers with
cost polynomial in triangulation size and logarithmic normal weights;
the isotopy result additionally assumes essential normal one-manifolds
\cite{lackenbycurves2026}. These provide concrete components for a geometric
implementation. Neither a fast terminal-pattern check nor an iteration
count by itself supplies a complete bound on construction, cutting,
simplification, representation growth and restarts.

For example, an independently proved bound $L=O(\log n)$ and
$G=n^{O(1)}$ would make $L(G+1)^L=n^{O(\log n)}$; with
$L=O((\log n)^2)$ it would give $n^{O((\log n)^2)}$.
These are conditional substitutions, not bounds established for the
hierarchy inputs generated by this package. Likewise, a bit cost polynomial
in an expanded normal weight is inadequate when that weight is exponential
in its encoded length.

The strongest credible program therefore has two interacting branches:
exact algebraic obstructions and compressed homological decisions for
efficient practical rejection, and fully costed compressed geometric
operations for a possible general recognition guarantee. They should share
validated inputs, deterministic evidence records and resource-safe
orchestration, while retaining separate mathematical obligations.
